\section{Đọc được không có nghĩa là điều khiển được}
\label{sec:ch15-doc-duoc-dieu-khien-duoc}

\index{sparse autoencoder!nút thắt không ai gán nhãn|(}%
Nhánh không giám sát mà mục trước bỏ lại có một chỗ khó mà cả bốn mục trên
không chạm tới: ở đó không có nhãn khái niệm, nên mốc của
định nghĩa~
ef{def:ch15-hai-diem-so} không tồn tại và không đo rò rỉ được. Mục này đọc bài thứ ba của
chương, bài đứng đúng ở nhánh ấy, và nó không đo \tn{rò rỉ}{leakage}. Nó đo hai thứ khác, và
và hai thứ ấy không đi cùng nhau.

\index{sparse autoencoder!vì sao nó thuộc chương này}%
Sparse autoencoder, gọi tắt là SAE theo phụ lục~\ref{app:viet-tat}, học dựng lại
kích hoạt của một mô hình đã huấn luyện xong bằng một biểu diễn thưa và nhiều
chiều hơn chính kích hoạt ấy. Mục~\ref{sec:ch07-attribution-thanh-phan} đã đặt
nó ở đâu trong bản đồ: nó không chấm điểm cho thành phần có sẵn mà học ra thành
phần mới, với hy vọng rằng thành phần mới thì mỗi cái mang đúng một nghĩa thay
vì nhiều nghĩa chồng lên nhau. Đặt cạnh CBM thì thấy hai bên là hai nửa của cùng
một ý. CBM có tên trước rồi mới có đơn vị: con người chọn tập khái niệm, và mô
hình phải học sao cho mỗi đơn vị khớp một cái tên. SAE có đơn vị trước rồi mới
đi tìm tên: nó học ra hàng vạn đơn vị thưa, và việc gán nghĩa cho từng đơn vị là
một bước riêng sau đó. Nên bài này là nút thắt khái niệm dựng theo chiều ngược,
và chỗ hỏng của nó cũng ngược lại.

\index{sparse autoencoder!hai phép đo trên từng nơ-ron}%
Bài báo đo hai tính chất trên từng nơ-ron của một SAE, và định nghĩa cả hai bằng
những phép tính rẻ~\autocite{p32cbmsae}. Tính chất thứ nhất là đọc được: một
công cụ dò tên có sẵn nhận vào một tập khái niệm bằng chữ và gán cho mỗi nơ-ron
cái tên khớp nhất với cách nó phản ứng trên một tập ảnh dò, còn điểm số là độ
khớp ấy. Tính chất thứ hai là \tn{khả năng điều khiển}{steerability}: ép kích
hoạt của đúng một nơ-ron lên một giá trị lớn, cho cả biểu diễn đã sửa chạy qua
bộ giải mã rồi qua một mô hình ngôn ngữ thị giác, hỏi nó bức ảnh này có gì bằng
đúng một từ, rồi đo độ gần giữa từ nó trả về và cái tên đã gán cho nơ-ron ấy.
Phép ép ấy có hai biến thể khác nhau ở chỗ những nơ-ron còn lại nhận giá trị gì
trong lúc ép, và bài báo chạy cả hai: hoặc đặt tất cả về không, hoặc để chúng ở
giá trị mà một bức ảnh trắng sinh ra. Cộng với hai cách chấm điểm đọc được, độ
khớp của công cụ dò tên và một điểm số đo mức đơn nghĩa của nơ-ron, thành bốn
chỉ số mà mọi bảng kết quả dưới đây báo cáo.

\index{mức quan trọng!đọc được và điều khiển được đo tách nhau}%
Hai tính chất ấy đo hai thứ mà Phần~IV đã phải tách ra bằng lời. Đọc được là một
phát biểu về những gì quan sát thấy: nơ-ron này sáng lên ở những ảnh nào, và cái
gì chung giữa chúng. Điều khiển được là một phát biểu về nhân quả: ép nơ-ron này
thì đầu ra đổi theo hướng nào. Và phát biểu thứ hai chính là mức quan trọng theo
nghĩa mà mục~\ref{sec:ch09-doi-tuong-khac-nhau} đã cố định, không phải độ trung
thực: nó chứng thực rằng đơn vị ấy có ảnh hưởng nhân quả lên đầu ra, không chứng
thực rằng cái tên gán cho nó mô tả đúng vai trò của nó.

\index{sparse autoencoder!bốn góc phần tư}%
Kết quả là một bảng bốn ô. Trên một SAE 65\,536 nơ-ron huấn luyện từ kích hoạt
của tầng 22 của một bộ mã hóa ảnh CLIP, với hệ số nở 64 và dữ liệu ImageNet-1k,
lấy trung bình của mỗi điểm số làm ranh giới, số nơ-ron rơi vào bốn góc là:
23\,763 nơ-ron, tức 36.26 phần trăm, thấp cả hai; 13\,022 nơ-ron, tức 19.87 phần
trăm, đọc được mà không điều khiển được; 16\,403 nơ-ron, tức 25.03 phần trăm,
điều khiển được mà không đọc được; và 12\,348 nơ-ron, tức 18.84 phần trăm, cao
cả hai. Bốn số ấy cộng lại đúng bằng 65\,536.

Con số cuối là con số cả mục xoay quanh: chưa tới một phần năm số nơ-ron của
một SAE vừa có tên đọc được vừa ép được đầu ra đi theo cái tên ấy. Hai góc giữa
nói vì sao con số ấy nhỏ. Có nơ-ron sáng lên đều đặn ở một loại ảnh, nhận được
một cái tên rõ ràng, và ép nó thì đầu ra không đổi theo hướng nào cả. Có nơ-ron
ép thì đầu ra đổi mạnh, mà tập ảnh làm nó sáng thì không có cái tên nào gọn cho
chúng. Tổng hai góc ấy là 44.90 phần trăm, gần một nửa từ điển.

\index{sparse autoencoder!khái niệm người dùng muốn thì không có}%
Chỗ hỏng thứ hai là chỗ hỏng về diện, và nó không phụ thuộc vào ranh giới trung
bình như bảng bốn ô. Dùng chính công cụ dò tên ấy để đếm xem trong một tập khái
niệm cho trước, bao nhiêu khái niệm được ít nhất một nơ-ron của SAE nhận, thì
kết quả phụ thuộc mạnh vào chỗ tập ấy lấy từ đâu. Bảng~\ref{tab:ch15-dien-khai-niem}
là năm tập mà bài báo thử.

\begin{table}[htbp]
  \centering
  \caption{Số khái niệm mà một SAE 65\,536 nơ-ron nhận được, trên năm tập khái
  niệm khác nhau~\autocite{p32cbmsae}. Hai dòng giữa là hai tập gắn với chính
  bộ dữ liệu mà SAE được huấn luyện trên, nên chúng là hai dòng khó bào chữa
  nhất. Hai dòng cuối cho thấy diện phủ đi xuống theo cỡ của tập chứ không theo
  mức độ tập ấy xa dữ liệu huấn luyện.}
  \label{tab:ch15-dien-khai-niem}
  \begin{tabular}{@{}lrrr@{}}
    \toprule
    Tập khái niệm & Nhận được & Tổng & Tỉ lệ \\
    \midrule
    Tập thị giác nhỏ dùng để chấm điểm đọc được & 1\,153 & 1\,197 & 96.3\% \\
    \addlinespace
    Tập thứ nhất gắn với ImageNet & 3\,445 & 4\,729 & 72.8\% \\
    Tập thứ hai gắn với ImageNet & 4\,333 & 7\,827 & 55.3\% \\
    \addlinespace
    Ba nghìn từ tiếng Anh thông dụng & 1\,857 & 3\,000 & 61.9\% \\
    Hai mươi nghìn từ tiếng Anh thông dụng & 5\,596 & 20\,000 & 28.0\% \\
    \bottomrule
  \end{tabular}
\end{table}

Dòng đầu là dòng dễ, và nó dễ vì tập ấy nhỏ và định nghĩa gọn. Hai dòng giữa là
chỗ đắt: chúng gắn với chính bộ dữ liệu mà SAE được huấn luyện trên, nên không
giải thích được bằng việc hỏi sai chỗ, và ở đó SAE thiếu 27 tới 45 phần trăm số
khái niệm. Hai dòng cuối cho thấy cái quyết định không phải khoảng cách tới dữ
liệu huấn luyện mà là cỡ của tập: hỏi ba nghìn từ thì thiếu gần bốn phần mười,
hỏi hai mươi nghìn từ thì thiếu hơn bảy phần mười. Một từ điển 65\,536 đơn vị,
tức lớn gấp hơn ba lần danh sách hai mươi nghìn từ ấy, vẫn không mang nổi một
phần ba số khái niệm trong đó.

\index{CB-SAE!ghép hai nhánh}%
Phép sửa mà bài đề xuất ghép hai nhánh của mục~\ref{sec:ch15-khai-niem-trong-mo-hinh}
lại. Nó chấm điểm từng nơ-ron của SAE theo hai tính chất trên, cắt bỏ những
nơ-ron có tổng hai điểm số thấp nhất, rồi bù vào phần đã mất bằng một tầng thắt
khái niệm học riêng, gắn với một tập khái niệm do người dùng nêu. Tầng bù ấy học
theo ba mục tiêu: dựng lại được kích hoạt như trước khi cắt, kích hoạt đúng
khái niệm của mình, và ép được đầu ra theo khái niệm ấy. Kết quả gọi là CB-SAE.
Cách nó chọn tập khái niệm cho tầng bù nói rõ ý ấy: nó chỉ lấy những khái
niệm mà SAE giữ lại \emph{không} biểu diễn được, nên tầng thắt ở đây không thay
SAE mà lấp chỗ SAE thiếu.

\index{CB-SAE!hai tính chất vẫn không đi cùng nhau}%
Đo lại trên ba mô hình phía sau, hai mô hình ngôn ngữ thị giác và một mô hình
sinh ảnh, CB-SAE tốt hơn SAE gốc ở cả bốn chỉ số của cả ba mô hình. Nhưng bảng
tách riêng ba nhóm nơ-ron mới là chỗ nói được nhiều nhất, và nó nói lại đúng
điều mà bốn góc phần tư đã nói. Trên mô hình ngôn ngữ thị giác thứ nhất, nhóm
nơ-ron thắt khái niệm mới đạt điểm đọc được 0.323, cao hơn hẳn nhóm nơ-ron SAE
được giữ lại ở 0.238; nhưng chính nhóm ấy lại điều khiển kém hơn, 0.231 so với
0.263 theo cách ép thứ nhất và 0.219 so với 0.252 theo cách ép thứ hai. Đặt tên
sẵn cho một đơn vị làm nó đọc được hơn và không làm nó điều khiển được hơn. Hai
tính chất tách nhau cả trong SAE lẫn trong phép sửa dựng ra để hàn chúng lại.

\index{CB-SAE!một con số cộng chứ không phải trung bình}%
Một chỗ trong phần kết quả cần sửa trước khi trích. Bài báo báo cáo sáu mức tăng
tương đối của CB-SAE so với SAE gốc, mỗi mức là trung bình hai chỉ số của một
cặp tính chất và mô hình, và gọi cả sáu là trung bình. Dựng lại từ chính bảng
của bài thì năm mức đúng là trung bình, còn mức thứ tư, khả năng điều khiển trên
mô hình ngôn ngữ thị giác thứ hai, là tổng: hai chỉ số của cặp ấy tăng 7.26 và
6.78 phần trăm, nên trung bình của chúng là 7.02 còn tổng là 14.04, và con số
bài in cho ô ấy là 14.0. Nó đi tiếp vào phần tóm tắt, nơi mức tăng khả năng điều
khiển của cả bài được báo là 14.5 phần trăm, trung bình của ba mức 27.5, 14.0 và
2.1. Thay 14.0 bằng 7.0 thì con số ấy là 12.2 phần trăm. Mức tăng đọc được, 32.1
phần trăm, thì
đúng: nó là trung bình của 33.0, 29.0 và 34.3, và cả ba mức ấy đều dựng lại
được. Sách in con số đã sửa và giữ nguyên mọi thứ hạng, vì dấu của mọi phép so
không đổi: CB-SAE vẫn hơn SAE gốc ở cả bốn chỉ số của cả ba mô hình.

\index{sparse autoencoder!một công cụ đứng cả hai phía}%
Còn một chỗ mà bài không nói và sách nói. Phần giới hạn của bài ghi rằng hiệu
quả của phương pháp phụ thuộc vào việc công cụ dò tên gán tên có chính xác hay
không, và đó là một câu đúng nhưng chưa đủ. Công cụ ấy không đứng một phía mà
đứng cả hai: nó gán cái tên, điểm số đọc được \emph{là} độ khớp mà nó tính ra
khi gán, và điểm số điều khiển được thì đo độ gần giữa đầu ra sau khi ép và
chính cái tên nó đã gán. Nên một nơ-ron mà công cụ ấy gán nhầm tên sẽ bị chấm
thấp ở cả hai điểm số cùng lúc, và bốn góc phần tư ở trên không tách được một
nơ-ron thật sự vô nghĩa khỏi một nơ-ron có nghĩa mà công cụ dò không gọi được
tên. Đây là hình dạng quen: một dụng cụ chấm điểm cho một lời giải thích, dựng
trên một dụng cụ khác chưa ai kiểm, đúng cái vòng mà
chương~\ref{ch:tro-choi-attribution} đã đi.

\index{sparse autoencoder!nút thắt không ai gán nhãn|)}%
\index{khoảng trống nghiên cứu!cái chương này để lại}%
Gom cả chương lại thì có ba thứ đi tiếp và một thứ dừng ở đây. Thứ đi tiếp thứ
nhất là đường thứ tư dựng ground truth mà
mục~\ref{sec:ch15-dung-cu-duoc-kiem} đã đặt tên, cùng cái giá của nó, và
chương~\ref{ch:ground-truth-dung-san} nhận nó cạnh ba đường kia. Cái giá ấy có
hai nửa: phải có người gán nhãn khái niệm cho từng điểm dữ liệu, và nếu tập nhãn
ấy không đủ dự báo nhiệm vụ thì theo chính lời người đề xuất khung ở
mục~\ref{sec:ch15-nguyen-nhan}, nhiệm vụ ấy không dùng lối này được. Thứ đi tiếp
thứ hai là một phép đo, và nó rẻ: dựng một tập khái niệm không cần nhãn thì làm
được, còn hỏi xem mỗi đơn vị có ép được đầu ra hay không thì chỉ tốn một lượt
chạy cho mỗi đơn vị. Thứ đi tiếp thứ ba là một yêu cầu đặt lên dụng cụ tương
lai, và nó là yêu cầu thứ ba mà sách gom được: mọi phép đọc tên của một đơn vị
bên trong mô hình phải đi kèm một phép thử ép đơn vị ấy, và phải báo cả hai kết
quả, vì 18.84 phần trăm là tỉ lệ những đơn vị mà hai kết quả cùng đạt. Yêu cầu
ấy đứng cạnh yêu cầu công bố tập can thiệp của
mục~\ref{sec:ch14-dieu-chua-giai-quyet} và yêu cầu mô tả chỗ lệch của
mục~\ref{sec:ch13-tran-noi-gi}.

Thứ dừng ở đây là chính chữ trong nhan đề sách. Không bài nào trong ba bài dùng
chữ \enquote{faithful} một lần nào trong phần thân, và điều đó lặp lại đúng cái
đã xảy ra ở chương~\ref{ch:giai-thich-duoi-tan-cong}: hai tầng liền của thang
đọc, sáu bài báo, không tầng nào với tới chữ ấy. Ba đại lượng mà chương này đọc
là lượng thông tin thừa trong một cột khái niệm, độ khớp giữa một cái tên và một
tập ảnh, và độ gần giữa một đầu ra bị ép và cái tên ấy. Không đại lượng nào so
một lời giải thích với hành vi của mô hình dưới một tập can thiệp đã nêu, tức
không đại lượng nào là định nghĩa~\ref{def:ch01-do-trung-thuc}.
Chương~\ref{ch:ground-truth-dung-san} đọc ba lối thoát mà văn liệu phê phán chỉ
ra, và câu hỏi nó mang theo từ đây là câu hỏi mà mục này vừa đo được một nửa:
một khái niệm bên trong mô hình đọc được tên, có nghĩa là gì.
